\section{Related work}\label{sec:2}\label{related-work}

\textbf{Finite-population inference.} Waudby-Smith and Ramdas provide bounded and categorical confidence sequences for sampling without replacement \citep{r1}. Their composite-null inversion already permits decisions about regions of the mean. Spertus, Sridhar, and Stark study sequential stratified inference, including adaptive interleaving and stronger union-intersection constructions \citep{r4}.

\textbf{Sequential model comparison.} Arviv et al.~combine group-sequential testing with pairwise comparison and equivalence margins \citep{r3}. Their stated calibration uses approximate normality, independent observations, and planned looks. We instead condition on the realized finite score vector and use non-asymptotic design-based guarantees. SySRs exploits paired observations and model similarity in fixed-budget best-model identification \citep{r5}. PULSE combines finite-pool inference, adaptive model allocation, and low-rank predictions \citep{r6}. Their objectives differ from certifying a three-region decision; fixed-budget accuracy cannot be compared to a stopping guarantee without an explicit adaptation.

\textbf{Prediction-assisted and adaptive evaluation.} CELEUS supplies anytime-valid finite-pool inference from surrogate-completed, inverse-probability-corrected signals \citep{r7}. It is a close general framework in which paired differences can be treated as a bounded outcome. Hsu and Shekhar study active confidence-sequence construction using historical model responses, with a target defined through per-question success probabilities \citep{r8}. Their analysis already highlights cases where uniform querying competes favorably with adaptive querying. Bayesian coreset sizing, including multi-target modeling, is developed by BayesAME \citep{r9}.

\textbf{Measurement and uncertainty.} ATLAS uses IRT and Fisher information for adaptive ability estimation \citep{r10}; Jiang et al.~examine reliability of IRT estimation in AI benchmark regimes \citep{r11}. Balkır et al.~extend adaptive ranking to continuous scores and cost-aware pairwise stopping \citep{r12}. Neuhof and Benjamini emphasize subject heterogeneity and the importance of selecting the inferential population \citep{r13}. We do not extrapolate fixed-matrix results to new subjects. Pilditch studies Bayesian precision and stability stopping, including repeated attempts \citep{r14}. AV-AIVAT combines variance reduction and confidence sequences for agent evaluation in games \citep{r15}. tinyBenchmarks is an important earlier example of efficient item selection and score reconstruction \citep{r16}.

\textbf{Sequential auditing beyond model evaluation.} Kato and Nakagawa \citep{r18} construct finite-population binary audit tests with hypergeometric recursion, prespecified looks and controlled errors outside an indifference region. Unlike affirmative equivalence here, their indifference region imposes no decision-error restriction. Their review also traces hypergeometric sequential testing to Lai (1979) and subsequent closed multistage procedures. RiLACS \citep{r19} uses confidence sequences to certify election assertions, continuing to a full count when necessary. Its mean-threshold reduction is closely related to each directional certificate used here. Risk-limiting financial audits \citep{r20} already combine weighted WoR sampling, bounded-value confidence sequences and side information. Thus fixed-population certification, full-count fallback and prediction-assisted sampling have substantial antecedents outside LLM evaluation.

\textbf{Closest cost and lower-bound comparisons.} Arviv et al.~already plot evaluation cost against performance difference; studying difficulty-dependent cost is therefore not itself new. Their current author code (commit \texttt{2c44ed23d2a8df2b42a036c9dc639bb71c0c07d5}) implements efficacy by exclusion of zero and equivalence by containment in the margin. Those conditions can both hold. Our mutually exclusive practical-superiority labels instead use the outer thresholds \(\pm\)\(\epsilon\). A numerical head-to-head using the original stop reasons would consequently score different claims. The main comparisons therefore use inference constructions with the same three labels. Harmonizing an original pipeline is possible, but requires changing its stopping target and checking its calibration; target mismatch is a reason not to compare raw stop reasons, not evidence that our evaluator is cheaper.

Shekhar and Ramdas \citep{r17} derive instance-dependent lower bounds for fixed-sample WoR CI widths in a moderate-accuracy regime. Their Theorem 3.4 and discussion of time-uniform extensions are directly relevant prior art. Our coupling result concerns adaptive query cost for a label-changing finite-vector alternative, including fixed-\(\alpha\) endpoint cases; it does not replace their richer width lower bounds or establish a general optimal rate. The edit-capacity argument specializes a standard two-instance indistinguishability principle to query cost. Its role is to expose which unseen changes obstruct this decision, rather than introduce a new general lower-bound method.
